\documentclass[11pt, a4paper]{article}

\usepackage[utf8]{inputenc}
\usepackage[T1]{fontenc}
\usepackage{amsmath, amssymb, amsthm, mathtools}
\usepackage{bm}
% indicator function via boldsymbol (no extra package needed)
\usepackage{booktabs}
\usepackage{graphicx}
\usepackage{hyperref}
\usepackage[margin=2.5cm]{geometry}
\usepackage{enumitem}
\usepackage{xcolor}
\usepackage{tcolorbox}
\usepackage{float}

% Theorem environments
\theoremstyle{definition}
\newtheorem{definition}{Definition}[section]
\newtheorem{example}{Example}[section]
\theoremstyle{plain}
\newtheorem{theorem}{Theorem}[section]
\newtheorem{proposition}{Proposition}[section]
\newtheorem{corollary}{Corollary}[section]
\theoremstyle{remark}
\newtheorem{remark}{Remark}[section]

% Colored boxes
\newtcolorbox{keyresult}[1][]{
    colback=blue!5!white,
    colframe=blue!60!black,
    fonttitle=\bfseries,
    title={Key Insight},
    #1
}
\newtcolorbox{warningbox}[1][]{
    colback=red!5!white,
    colframe=red!60!black,
    fonttitle=\bfseries,
    title={Warning},
    #1
}
\newtcolorbox{refbox}[1][]{
    colback=gray!5!white,
    colframe=gray!50!black,
    fonttitle=\bfseries,
    title={Go Deeper},
    #1
}
\newtcolorbox{keyidea}[1][]{
    colback=green!5!white,
    colframe=green!50!black,
    fonttitle=\bfseries,
    title={Core Idea},
    #1
}
\newtcolorbox{prereq}[1][]{
    colback=orange!5!white,
    colframe=orange!60!black,
    fonttitle=\bfseries,
    title={Read Before Coding},
    #1
}
\newtcolorbox{objective}[1][]{
    colback=blue!5!white,
    colframe=blue!60!black,
    fonttitle=\bfseries,
    title={What You Will Build},
    #1
}
\newtcolorbox{interpbox}[1][]{
    colback=green!5!white,
    colframe=green!50!black,
    fonttitle=\bfseries,
    title={Interpretation},
    #1
}

% Custom commands
\newcommand{\VaR}{\operatorname{VaR}}
\newcommand{\ES}{\operatorname{ES}}
\newcommand{\CVaR}{\operatorname{CVaR}}
\newcommand{\EV}{\mathbb{E}}
\newcommand{\PV}{\mathbb{P}}
\newcommand{\RV}{\mathbb{R}}
\newcommand{\ind}{\boldsymbol{1}}
\newcommand{\PnL}{\operatorname{P\&L}}

\title{%
    \textbf{Stress Testing Framework} \\[0.5em]
    \large Measuring Portfolio Resilience Under Extreme Scenarios
}
\author{Andr\'es Gonz\'alez Ortega}
\date{March 2026}

\begin{document}
\maketitle

\begin{abstract}
An applied, example-driven treatment of stress testing for financial portfolios.
We begin from the practical limitations of statistical risk measures (VaR, ES)
that are calibrated to ``normal'' market conditions,
and develop the three pillars of stress testing:
historical scenario replay, hypothetical (regulatory) scenarios, and reverse stress testing.
The Fed DFAST architecture is examined in detail,
including the macro-to-portfolio transmission mechanism via regression-based and satellite models.
Credit migration under stress and reverse stress testing as an optimization problem are treated
with worked numerical examples throughout.
The document assumes familiarity with VaR/ES (Project~03), volatility modeling (Project~04),
and tail risk (Project~05).
\end{abstract}

\tableofcontents
\newpage

%% ============================================================================
\section{Why Stress Testing}
\label{sec:why-stress}
%% ============================================================================

\begin{example}[The 2008 wake-up call]\label{ex:2008-wakeup}
In the first half of 2007, the $\VaR_{0.99}$ for a diversified portfolio of US equities
and mortgage-backed securities was typically estimated at 2--3\% of portfolio value (daily).
The GARCH models were calibrated to 2003--2007 data --- a period of unusually low volatility
(the ``Great Moderation'').

Between September and November 2008, daily portfolio losses of 5--8\%
occurred on multiple occasions.
The 99\% VaR was breached not once or twice, but \emph{dozens} of times in a single quarter.
Models that were ``correct'' under recent history
were catastrophically wrong under crisis conditions.
\end{example}

The failure was not that VaR or ES were computed incorrectly.
The failure was \textbf{structural}: statistical risk measures answer the question
``what is the worst loss under the distribution estimated from recent data?''
They do \emph{not} answer the question
``what happens if the world changes in a specific, adverse way?''

Stress testing is the complement to statistical risk measures.
It asks a fundamentally different question:

\begin{keyidea}
\textbf{Statistical risk measures} (VaR, ES):
``Given the current regime, what is the $\alpha$-quantile or tail expectation of the loss distribution?''

\medskip

\textbf{Stress testing}:
``Given a specific adverse scenario (whether historical or hypothetical),
what would the portfolio lose?''

\medskip

The two approaches are complementary, not competing.
VaR/ES measures the tail \emph{within} the estimated distribution.
Stress testing evaluates losses \emph{outside} the estimated distribution ---
under scenarios that the model may assign negligible probability to,
but that regulators, risk managers, and history tell us can happen.
\end{keyidea}

The regulatory mandate is explicit and growing:

\begin{itemize}
    \item \textbf{Basel III/IV} (BCBS, 2018): banks must conduct regular stress tests
          covering credit, market, and operational risk.
          Internal stress testing is a pillar of the Internal Capital Adequacy Assessment Process (ICAAP).
    \item \textbf{DFAST / CCAR} (US): the Dodd-Frank Act Stress Tests require
          all bank holding companies with assets above \$100B to project losses
          under Fed-prescribed macroeconomic scenarios.
    \item \textbf{EBA / ECB} (EU): the European Banking Authority conducts
          biennial EU-wide stress tests with prescribed adverse scenarios.
    \item \textbf{PRA} (UK): the Bank of England's Prudential Regulation Authority
          runs concurrent stress tests for major UK banks.
\end{itemize}

\begin{warningbox}
The 2008 crisis did not just reveal that models were miscalibrated.
It revealed a \textbf{conceptual gap}: risk management frameworks built entirely on
statistical measures estimated from recent data will always be blind to regime shifts.
Stress testing fills this gap by forcing institutions to evaluate losses under scenarios
that are not in the historical sample --- or that occurred so long ago
that statistical models have ``forgotten'' them.
\end{warningbox}

%% ============================================================================
\section{Types of Stress Tests}
\label{sec:types}
%% ============================================================================

Stress tests come in three broad categories, each answering a different question.

\subsection{Historical Scenarios}

\textbf{Approach}: take the actual market moves observed during a past crisis
and apply them to the current portfolio.

\textbf{Question answered}: ``If the 2008 GFC happened again today, what would we lose?''

\textbf{Advantages}:
\begin{itemize}
    \item Scenarios are internally consistent (all factor moves actually co-occurred).
    \item Easy to communicate to senior management (``this is what happened last time'').
    \item No modeling assumptions needed --- just apply the historical factor changes.
\end{itemize}

\textbf{Limitations}:
\begin{itemize}
    \item The next crisis may look different from any past crisis.
    \item The portfolio composition today may have exposures to risk factors
          that did not exist or were not material in the historical period.
    \item Historical scenarios are backward-looking by definition.
\end{itemize}

\subsection{Hypothetical Scenarios}

\textbf{Approach}: construct a scenario from economic logic or regulatory specification,
specifying paths for key macroeconomic variables over a projection horizon.

\textbf{Question answered}: ``If GDP fell by 6\% and unemployment rose to 10\%,
what would our portfolio lose?''

\textbf{Examples}: the Fed DFAST baseline, adverse, and severely adverse scenarios;
internal scenarios designed by a bank's economics team.

\textbf{Advantages}:
\begin{itemize}
    \item Forward-looking: can incorporate risks that have not yet materialized.
    \item Tailored to specific vulnerabilities of the institution.
    \item Allows exploration of ``what if'' questions that history has not answered.
\end{itemize}

\textbf{Limitations}:
\begin{itemize}
    \item The scenario may be internally inconsistent
          (e.g., GDP falls but equity markets rise --- is that plausible?).
    \item Requires a transmission model from macro variables to portfolio losses.
    \item Results are only as good as the transmission model.
\end{itemize}

\subsection{Reverse Stress Testing}

\textbf{Approach}: instead of specifying a scenario and computing the loss,
specify a loss threshold and find the scenario that produces it.

\textbf{Question answered}: ``What scenario would cause our portfolio
to lose enough to breach our capital ratio?''

\textbf{Advantages}:
\begin{itemize}
    \item Identifies hidden vulnerabilities that forward stress tests may miss.
    \item Forces the institution to think about its ``break point.''
    \item Required by several regulators (PRA, EBA).
\end{itemize}

\textbf{Limitations}:
\begin{itemize}
    \item Computationally harder (inverse problem).
    \item Multiple scenarios may produce the same loss --- which one is ``the'' answer?
    \item Requires careful formulation of the optimization problem.
\end{itemize}

\subsection{Comparison}

\begin{table}[H]
\centering
\begin{tabular}{@{}p{3cm}p{3.5cm}p{3.5cm}p{3.5cm}@{}}
\toprule
& \textbf{Historical} & \textbf{Hypothetical} & \textbf{Reverse} \\
\midrule
\textbf{Input} & Past crisis dates & Macro variable paths & Loss threshold \\
\textbf{Output} & Portfolio loss & Portfolio loss & Scenario \\
\textbf{Direction} & Scenario $\to$ Loss & Scenario $\to$ Loss & Loss $\to$ Scenario \\
\textbf{Key strength} & Realism, consistency & Forward-looking & Finds blind spots \\
\textbf{Key weakness} & Backward-looking & Model-dependent & Computational cost \\
\textbf{Regulatory use} & Supplementary & DFAST/CCAR, EBA & PRA, ICAAP \\
\bottomrule
\end{tabular}
\caption{Comparison of the three stress testing approaches.}
\end{table}

\begin{keyresult}
A robust stress testing framework uses \emph{all three} approaches.
Historical scenarios provide a reality anchor.
Hypothetical scenarios explore forward-looking risks.
Reverse stress testing ensures no critical vulnerability is overlooked.
The three approaches are complementary, not substitutes.
\end{keyresult}

%% ============================================================================
\section{Fed DFAST Scenario Architecture}
\label{sec:dfast}
%% ============================================================================

The Federal Reserve's Dodd-Frank Act Stress Tests (DFAST) provide the most structured
example of hypothetical stress testing.
Each year, the Fed publishes three macroeconomic scenarios,
and covered institutions must project their capital ratios under each.

\subsection{The Three Scenarios}

\begin{enumerate}
    \item \textbf{Baseline}: a ``most likely'' economic path.
          Roughly consistent with the Blue Chip consensus forecast.
          GDP growth $\approx$ 2\%, unemployment $\approx$ 4\%, equity indices roughly flat.
    \item \textbf{Adverse}: a moderate recession.
          GDP growth $\approx$ $-1$\% to $-2$\%, unemployment rises to $\approx$ 7\%,
          equity indices decline $\approx$ 25\%, credit spreads widen moderately.
    \item \textbf{Severely adverse}: a deep recession.
          GDP growth $\approx$ $-6$\% to $-8$\%, unemployment rises to $\approx$ 10\%,
          equity indices decline $\approx$ 55\%, credit spreads spike,
          house prices fall $\approx$ 25\%.
\end{enumerate}

\subsection{Macro Variable Paths}

The Fed specifies quarterly paths for 28 macroeconomic and financial variables
over a 9-quarter (or 13-quarter) projection horizon.
The key variables include:

\begin{table}[H]
\centering
\begin{tabular}{@{}lcccc@{}}
\toprule
\textbf{Variable} & \textbf{Baseline} & \textbf{Adverse} & \textbf{Severely Adverse} & \textbf{Units} \\
\midrule
Real GDP growth (Y/Y) & $+2.0\%$ & $-2.0\%$ & $-6.0\%$ & \% \\
Unemployment rate (peak) & $4.2\%$ & $7.0\%$ & $10.0\%$ & \% \\
Dow Jones / S\&P trough & $-2\%$ & $-25\%$ & $-55\%$ & \% from current \\
10Y Treasury yield (trough) & $3.8\%$ & $2.5\%$ & $1.0\%$ & \% \\
BBB corporate spread (peak) & $+50$bps & $+200$bps & $+450$bps & bps \\
House Price Index (HPI) & $+2\%$ & $-10\%$ & $-25\%$ & \% \\
VIX (peak) & $20$ & $45$ & $70$ & index \\
\bottomrule
\end{tabular}
\caption{Representative DFAST macro variable paths (stylized from recent exercises).
Actual Fed scenarios vary year to year.}
\end{table}

\subsection{What Banks Must Do}

For each scenario, covered institutions must:
\begin{enumerate}
    \item Project quarterly \textbf{pre-provision net revenue} (PPNR).
    \item Project quarterly \textbf{credit losses} (provision for loan losses).
    \item Project \textbf{trading and counterparty losses} (for firms with large trading books).
    \item Project \textbf{operational risk losses}.
    \item Compute \textbf{capital ratios} (CET1, Tier 1, Total Capital, Leverage)
          at each quarter of the horizon.
    \item Demonstrate that capital ratios stay above regulatory minimums
          \emph{at every quarter}, not just at the end.
\end{enumerate}

\begin{warningbox}
The DFAST severely adverse scenario is designed to be harsh but \emph{plausible}.
It is not a ``worst case'' --- the 2008 crisis was actually more severe on some dimensions
(e.g., the S\&P~500 fell 57\% peak-to-trough, and unemployment reached 10\%).
The regulatory scenario is a floor, not a ceiling, for stress testing.
Institutions are expected to run their own, potentially more severe, internal scenarios.
\end{warningbox}

\subsection{Worked Example: Severely Adverse Impact}

\begin{example}[Severely adverse scenario]\label{ex:dfast-severe}
Consider a bank with a \$50 billion loan portfolio (60\% commercial, 40\% consumer)
and a \$10 billion trading book (equity-heavy).

Under the severely adverse scenario:
\begin{itemize}
    \item \textbf{Credit losses}: the loan portfolio's weighted average PD increases from 1.2\%
          (baseline) to 4.5\% (severely adverse), driven by GDP contraction and unemployment.
          Expected credit losses: $\$50\text{B} \times (4.5\% - 1.2\%) \times 0.45 \;\text{(LGD)} \approx \$743\text{M}$.
    \item \textbf{Trading losses}: equity index $-55\%$, credit spreads $+450$bps.
          Estimated mark-to-market loss on the trading book: $\approx \$2.8\text{B}$.
    \item \textbf{PPNR offset}: even under stress, the bank earns net interest income and fees.
          Projected 9-quarter cumulative PPNR: $\approx \$1.5\text{B}$.
    \item \textbf{Net stress loss}: $\$743\text{M} + \$2.8\text{B} - \$1.5\text{B} = \$2.04\text{B}$.
\end{itemize}

If the bank's CET1 capital is \$5.0 billion and risk-weighted assets are \$55 billion,
the pre-stress CET1 ratio is $5.0 / 55 = 9.1\%$.
After the stress loss: $(5.0 - 2.04) / 55 = 5.4\%$.
The regulatory minimum CET1 ratio is 4.5\%, plus a stress capital buffer.
The bank survives, but with minimal headroom.
\end{example}

\begin{interpbox}
The DFAST exercise is not just a compliance exercise.
It reveals the \textbf{capital adequacy} of the institution under severe but plausible conditions.
A bank that barely passes the severely adverse scenario has little buffer against
an even more severe event --- and the 2008 crisis was, on several dimensions,
more severe than the typical DFAST scenario.
\end{interpbox}

%% ============================================================================
\section{Macro-to-Portfolio Transmission}
\label{sec:transmission}
%% ============================================================================

The core modeling challenge in hypothetical stress testing:
how do macroeconomic shocks translate into portfolio losses?
The scenario gives us $\Delta\text{GDP}$, $\Delta\text{Unemployment}$, $\Delta\text{Equity Index}$, etc.
We need to convert these into a dollar loss for the portfolio.

Two main approaches exist.

\subsection{Approach 1: Regression-Based (Top-Down)}

\begin{keyidea}
\textbf{The regression approach}: regress portfolio returns (or P\&L) on macroeconomic factors
using historical data.
Then apply the stressed factor values to the estimated model to project losses.

\medskip

Let $R_t$ be the portfolio return at time $t$ and $\bm{X}_t = (X_{1,t}, \ldots, X_{k,t})^\top$
be a vector of macro factor changes. The model is:
\[
    R_t = \alpha + \sum_{i=1}^{k} \beta_i \, X_{i,t} + \varepsilon_t
\]

The \textbf{stressed P\&L} is:
\[
    \Delta \PnL_{\text{stress}} = V \cdot \left(\hat\alpha + \sum_{i=1}^{k} \hat\beta_i \, \Delta X_{i}^{\text{stress}}\right)
\]
where $V$ is the portfolio value and $\Delta X_i^{\text{stress}}$ is the change in macro factor $i$
under the stress scenario.
\end{keyidea}

\begin{example}[Regression-based transmission]\label{ex:regression}
Fit a quarterly regression of equity portfolio returns on macro factors
using 40 quarters (10 years) of data:
\[
    R_t = \alpha + \beta_1 \Delta\text{GDP}_t + \beta_2 \Delta\text{Unemp}_t
    + \beta_3 \Delta\text{Spread}_t + \varepsilon_t
\]

Estimated coefficients (representative):
\begin{center}
\begin{tabular}{@{}lccc@{}}
\toprule
\textbf{Factor} & $\hat\beta$ & \textbf{Stressed $\Delta X$} & $\hat\beta \times \Delta X$ \\
\midrule
$\Delta$GDP (Y/Y, \%) & $+2.8$ & $-6.0$ & $-16.8\%$ \\
$\Delta$Unemployment (\%) & $-3.5$ & $+5.8$ & $-20.3\%$ \\
$\Delta$BBB Spread (bps/100) & $-8.2$ & $+4.5$ & $-36.9\%$ \\
\midrule
Intercept $\hat\alpha$ & \multicolumn{2}{r}{} & $+1.2\%$ \\
\midrule
\textbf{Total stressed return} & \multicolumn{2}{r}{} & $\mathbf{-72.8\%}$ \\
\bottomrule
\end{tabular}
\end{center}

\textbf{Problem}: the linear regression projects a $-72.8\%$ loss,
which exceeds the 2008 actual equity decline of $-57\%$.
The model extrapolates linearly beyond the range of observed data,
and the three factor effects may overlap (multicollinearity).

\textbf{Mitigation}: cap the projected loss at a historically calibrated maximum,
use regularized regression (Ridge, Lasso), or use the satellite approach instead.
\end{example}

\begin{warningbox}
The regression approach is simple but fragile.
Key pitfalls:
\begin{itemize}
    \item \textbf{Multicollinearity}: GDP, unemployment, and credit spreads are highly correlated.
          OLS coefficients become unstable and can have wrong signs.
    \item \textbf{Linearity}: the true relationship between macro shocks and portfolio losses
          is almost certainly nonlinear, especially under extreme stress.
    \item \textbf{Extrapolation}: the stress scenario is, by design,
          outside the range of historical data used to estimate the model.
          Linear extrapolation can produce absurd results.
    \item \textbf{Structural breaks}: the sensitivity of the portfolio to macro factors
          may change over time as portfolio composition evolves.
\end{itemize}
\end{warningbox}

\subsection{Approach 2: Satellite Models (Bottom-Up)}

\begin{keyidea}
Instead of one regression for the entire portfolio,
build \textbf{separate models} for each component of portfolio loss,
each driven by the relevant macro factors:
\begin{itemize}
    \item \textbf{PD model}: $\text{PD}_t = f_1(\text{GDP}_t, \text{Unemp}_t, \ldots)$
    \item \textbf{LGD model}: $\text{LGD}_t = f_2(\text{HPI}_t, \text{Spread}_t, \ldots)$
    \item \textbf{Prepayment model}: $\text{CPR}_t = f_3(\text{Rate}_t, \text{HPI}_t, \ldots)$
    \item \textbf{Revenue model}: $\text{PPNR}_t = f_4(\text{Rate}_t, \text{GDP}_t, \ldots)$
\end{itemize}
Each satellite model is estimated independently, using data and features
appropriate to the specific risk component.
The stressed loss is assembled by feeding macro scenario paths into each model
and aggregating.
\end{keyidea}

The satellite approach is what large banks actually use for DFAST/CCAR submissions.
It is more granular, more defensible, and more transparent than the top-down regression.

\begin{example}[Satellite PD model]\label{ex:satellite-pd}
A logistic regression for commercial loan default probability:
\[
    \ln\!\left(\frac{\text{PD}_t}{1 - \text{PD}_t}\right)
    = \gamma_0 + \gamma_1 \,\text{GDP}_t + \gamma_2 \,\text{Unemp}_t + \gamma_3 \,\text{Spread}_t
\]
Estimated on 80 quarters of data (20 years), with quarterly PD computed
from the observed default rate on the bank's commercial loan portfolio.

Representative estimates: $\hat\gamma_0 = -4.2$, $\hat\gamma_1 = -0.15$,
$\hat\gamma_2 = 0.22$, $\hat\gamma_3 = 0.35$.

Under the severely adverse scenario (GDP $= -6\%$, Unemp $= 10\%$, Spread $= +450$bps):
\[
    \text{logit}(\text{PD}_{\text{stress}}) = -4.2 + (-0.15)(-6) + 0.22(10) + 0.35(4.5) = 0.075
\]
\[
    \text{PD}_{\text{stress}} = \frac{1}{1 + e^{-0.075}} \approx 51.9\%?
\]

This is clearly too high --- the issue is that the logistic model is extrapolating
well beyond the training range.
In practice, banks constrain PD projections using expert judgment,
historical maximums, and model overlays.
A more realistic stressed PD for a commercial loan portfolio would be 4--6\%,
achieved by calibrating the model with additional controls.
\end{example}

\begin{keyresult}
Both the regression and satellite approaches face the same fundamental challenge:
stress scenarios are, \textbf{by construction}, outside the range of historical data.
Any model --- linear regression, logistic regression, machine learning ---
must extrapolate beyond its training domain.
This is why stress testing always requires \textbf{expert judgment} as a complement to models.
The model provides a disciplined starting point; judgment provides the guardrails.
\end{keyresult}

%% ============================================================================
\section{Credit Migration Under Stress}
\label{sec:migration}
%% ============================================================================

For credit portfolios, losses arise not only from outright defaults
but from \textbf{rating downgrades} that reduce the market value of bonds and loans.
A transition matrix describes how borrowers migrate between rating categories over time.

\subsection{The Transition Matrix}

\begin{definition}[Credit transition matrix]
A $K \times K$ \textbf{transition matrix} $\bm{P} = [p_{ij}]$ where $p_{ij}$ is the probability
that a borrower rated $i$ at time $t$ is rated $j$ at time $t+1$.
Rows sum to one: $\sum_{j=1}^{K} p_{ij} = 1$ for all $i$.
The last column ($j = K$) represents default, which is an absorbing state.
\end{definition}

\begin{example}[One-year transition matrix (stylized)]\label{ex:transition}
A simplified 4-state matrix (A, BBB, BB, Default):

\begin{center}
\begin{tabular}{@{}lcccc@{}}
\toprule
& \textbf{A} & \textbf{BBB} & \textbf{BB} & \textbf{Default} \\
\midrule
\textbf{A}   & 91.5\% & 7.0\% & 1.2\% & 0.3\% \\
\textbf{BBB} & 4.0\%  & 88.0\% & 6.5\% & 1.5\% \\
\textbf{BB}  & 0.5\%  & 6.0\%  & 85.0\% & 8.5\% \\
\textbf{Default} & 0\% & 0\% & 0\% & 100\% \\
\bottomrule
\end{tabular}
\end{center}

Under normal conditions, an A-rated borrower has a 91.5\% chance of remaining A,
a 7.0\% chance of being downgraded to BBB, and only a 0.3\% chance of defaulting.
\end{example}

\subsection{Stressing the Transition Matrix}

Under economic stress, two things happen simultaneously:
\begin{enumerate}
    \item \textbf{Downgrades accelerate}: the probability of moving to a lower rating increases.
    \item \textbf{Upgrades slow}: the probability of moving to a higher rating decreases.
\end{enumerate}

Two approaches to stress the matrix.

\subsubsection{Approach 1: Multiplicative Stress Factors}

Multiply the off-diagonal downgrade probabilities by a stress factor $s > 1$
and the off-diagonal upgrade probabilities by a factor $1/s$,
then rescale rows to sum to one.

For a stress factor $s = 2.0$ applied to Example~\ref{ex:transition}:

\begin{center}
\begin{tabular}{@{}lcccc@{}}
\toprule
& \textbf{A} & \textbf{BBB} & \textbf{BB} & \textbf{Default} \\
\midrule
\textbf{A}   & 83.0\% & 14.0\% & 2.4\% & 0.6\% \\
\textbf{BBB} & 2.0\%  & 82.0\% & 13.0\% & 3.0\% \\
\textbf{BB}  & 0.25\% & 3.0\%  & 79.75\% & 17.0\% \\
\bottomrule
\end{tabular}
\end{center}

The A-rated default probability doubled from 0.3\% to 0.6\%;
the BB default probability doubled from 8.5\% to 17.0\%.

\subsubsection{Approach 2: Generator Matrix (Matrix Exponential)}

\begin{keyidea}
A more rigorous approach uses the \textbf{generator matrix} $\bm{Q}$ such that
$\bm{P} = \exp(\bm{Q})$ (matrix exponential).
To stress the matrix, scale the generator:
\[
    \bm{P}_{\text{stress}} = \exp(s \cdot \bm{Q}), \quad s > 1
\]
This is equivalent to ``speeding up'' the credit deterioration process
by a factor of $s$.
The resulting stressed matrix is guaranteed to be a valid transition matrix
(all entries non-negative, rows sum to one) --- unlike the multiplicative approach,
which requires ad-hoc rescaling.

\medskip

Formally, $\bm{Q} = \ln(\bm{P})$ is the matrix logarithm (when it exists with non-negative off-diagonal entries).
The generator has the property that $q_{ij} \geq 0$ for $i \neq j$
and $q_{ii} = -\sum_{j \neq i} q_{ij}$ (rows sum to zero).
\end{keyidea}

\subsection{Worked Example: Investment-Grade Portfolio Under Stress}

\begin{example}[Stressed migration losses]\label{ex:stressed-migration}
Consider a \$10 billion investment-grade bond portfolio:
60\% rated A, 40\% rated BBB.
Under the stressed transition matrix ($s = 2.0$):

\textbf{Rating distribution after one year}:
\begin{itemize}
    \item Originally A (\$6B):
    $83.0\%$ stays A, $14.0\%$ becomes BBB, $2.4\%$ becomes BB, $0.6\%$ defaults.
    \item Originally BBB (\$4B):
    $2.0\%$ upgrades to A, $82.0\%$ stays BBB, $13.0\%$ becomes BB, $3.0\%$ defaults.
\end{itemize}

\textbf{Default losses} (assuming LGD = 60\%):
\[
    \text{Default loss} = (6 \times 0.006 + 4 \times 0.030) \times 0.60 = (0.036 + 0.120) \times 0.60 = \$93.6\text{M}
\]

\textbf{Mark-to-market migration losses}:
when a bond is downgraded, its spread widens and its market value falls.
Assume the following spread impacts:
\begin{itemize}
    \item A $\to$ BBB: spread widens by 50bps $\Rightarrow$ price drop $\approx$ 2.5\% (at 5Y duration).
    \item A $\to$ BB: spread widens by 200bps $\Rightarrow$ price drop $\approx$ 10\%.
    \item BBB $\to$ BB: spread widens by 150bps $\Rightarrow$ price drop $\approx$ 7.5\%.
\end{itemize}

Migration losses:
\begin{align*}
    \text{A} \to \text{BBB}: &\quad 6 \times 0.14 \times 0.025 = \$21.0\text{M} \\
    \text{A} \to \text{BB}: &\quad 6 \times 0.024 \times 0.10 = \$14.4\text{M} \\
    \text{BBB} \to \text{BB}: &\quad 4 \times 0.13 \times 0.075 = \$39.0\text{M}
\end{align*}

\textbf{Total stressed loss}: $\$93.6\text{M} + \$21.0\text{M} + \$14.4\text{M} + \$39.0\text{M} = \$168.0\text{M}$
(1.68\% of portfolio value).

Under the unstressed matrix, the total loss would be approximately \$60M.
Stress nearly triples the loss.
\end{example}

\begin{interpbox}
The majority of the stressed loss comes not from defaults but from \textbf{downgrades}.
This is typical of investment-grade portfolios: defaults are rare even under stress,
but the probability mass shifts downward through the rating spectrum,
causing widespread spread widening and mark-to-market losses.
For high-yield portfolios, the balance shifts toward default losses.
\end{interpbox}

%% ============================================================================
\section{Reverse Stress Testing}
\label{sec:reverse}
%% ============================================================================

Reverse stress testing inverts the standard question.
Instead of ``given a scenario, what is the loss?''
it asks ``given a loss threshold, what is the scenario?''

\subsection{The Optimization Formulation}

\begin{definition}[Reverse stress test]
Let $\bm{\Delta X} = (\Delta X_1, \ldots, \Delta X_k)^\top$ be a vector of macro factor changes
and $L(\bm{\Delta X})$ be the portfolio loss as a function of those changes.
Let $L^*$ be the loss threshold (e.g., the loss that would breach the minimum capital ratio).

The reverse stress test solves:
\[
    \min_{\bm{\Delta X}} \; \|\bm{\Delta X}\|_{\bm{\Sigma}^{-1}}
    \quad \text{subject to} \quad L(\bm{\Delta X}) \geq L^*
\]
where $\|\bm{\Delta X}\|_{\bm{\Sigma}^{-1}} = \sqrt{\bm{\Delta X}^\top \bm{\Sigma}^{-1} \bm{\Delta X}}$
is the Mahalanobis norm using the historical covariance matrix $\bm{\Sigma}$ of macro factor changes.
\end{definition}

\begin{keyidea}
The Mahalanobis norm measures the ``plausibility'' of the scenario.
A scenario where GDP falls by 10\% \emph{and} unemployment falls by 5\%
(contradictory) has a large Mahalanobis distance because such a combination
has never been observed.
A scenario where GDP falls by 6\% and unemployment rises to 10\%
(consistent with historical recessions) has a smaller Mahalanobis distance.

\medskip

Minimizing $\|\bm{\Delta X}\|_{\bm{\Sigma}^{-1}}$ subject to $L \geq L^*$
finds the \textbf{most plausible scenario} (closest to the mean, in a covariance-adjusted sense)
that produces the target loss.
This is the ``least surprising'' way to reach the break point.
\end{keyidea}

\subsection{Solution Approaches}

\subsubsection{Gradient-Based Optimization}

If $L(\bm{\Delta X})$ is differentiable (e.g., when using the linear regression model from Section~\ref{sec:transmission}),
the problem becomes:
\[
    \min_{\bm{\Delta X}} \; \bm{\Delta X}^\top \bm{\Sigma}^{-1} \bm{\Delta X}
    \quad \text{s.t.} \quad \bm{\beta}^\top \bm{\Delta X} \leq -L^* / V
\]
This is a quadratic program with a linear constraint --- solvable in closed form via the KKT conditions:
\[
    \bm{\Delta X}^* = -\frac{L^* / V}{\bm{\beta}^\top \bm{\Sigma} \bm{\beta}} \, \bm{\Sigma} \bm{\beta}
\]

The optimal scenario moves each factor in proportion to its beta times its covariance with other factors.
Factors with high betas \emph{and} high variance dominate the scenario.

\subsubsection{Grid Search}

When $L(\bm{\Delta X})$ is nonlinear (satellite models, option-like payoffs):
\begin{enumerate}
    \item Define a grid of plausible factor changes for each macro variable.
    \item Evaluate $L(\bm{\Delta X})$ at each grid point.
    \item Among all grid points where $L \geq L^*$,
          select the one with the smallest Mahalanobis distance.
\end{enumerate}

For $k = 5$ factors with 10 grid points each, this requires $10^5 = 100{,}000$ evaluations ---
feasible for fast loss functions, prohibitive for slow simulation-based models.
Latin hypercube sampling or importance sampling can reduce the computational cost.

\subsection{Worked Example: Finding the Break Point}

\begin{example}[Reverse stress test]\label{ex:reverse}
A bank has CET1 capital of \$5.0B, risk-weighted assets of \$55B,
and a minimum CET1 ratio requirement of 4.5\% (i.e., minimum capital of \$2.475B).
Projected PPNR over the stress horizon: \$1.5B.

The loss threshold that would breach the minimum:
\[
    L^* = \text{Capital} + \text{PPNR} - \text{Minimum Capital}
    = 5.0 + 1.5 - 2.475 = \$4.025\text{B}
\]

Using the linear model with $V = \$60\text{B}$, $\bm{\beta} = (2.8, -3.5, -8.2)^\top$
(from Example~\ref{ex:regression}), and the historical covariance matrix $\bm{\Sigma}$:

\begin{center}
\begin{tabular}{@{}lccc@{}}
\toprule
& $\Delta$GDP & $\Delta$Unemp & $\Delta$Spread \\
\midrule
$\Delta$GDP & $4.0$ & $-2.5$ & $-1.2$ \\
$\Delta$Unemp & $-2.5$ & $3.0$ & $1.8$ \\
$\Delta$Spread & $-1.2$ & $1.8$ & $2.5$ \\
\bottomrule
\end{tabular}
\end{center}

Solving the KKT conditions yields the critical scenario:
$\Delta\text{GDP}^* \approx -4.1\%$, $\Delta\text{Unemp}^* \approx +4.2\%$,
$\Delta\text{Spread}^* \approx +310\text{bps}$.

This is a scenario less severe than the DFAST severely adverse case,
which means the bank's capital buffer above the minimum is thin ---
a moderate recession, not a deep one, would be sufficient to breach the floor.
\end{example}

\begin{warningbox}
Reverse stress testing often reveals uncomfortable truths.
If the ``break point'' scenario is a moderate recession rather than an extreme crisis,
the institution's capital position is weaker than the forward stress tests alone might suggest.
This is precisely why reverse stress testing is valuable ---
and why some institutions resist implementing it rigorously.
\end{warningbox}

%% ============================================================================
\section{Historical Scenario Replay}
\label{sec:historical}
%% ============================================================================

Historical scenario replay is the most intuitive form of stress testing:
take the factor moves that occurred during a past crisis and apply them to the current portfolio.

\subsection{The Method}

\begin{enumerate}
    \item \textbf{Select a crisis period} and define the stress window
          (e.g., September 15, 2008 to March 9, 2009 for the GFC).
    \item \textbf{Extract factor changes}: compute the cumulative change in each risk factor
          over the stress window.
    \item \textbf{Revalue the portfolio}: apply the stressed factor changes
          to the current portfolio positions using either full revaluation
          or sensitivity-based (delta-gamma) approximation.
    \item \textbf{Compute the loss}: the difference between current portfolio value
          and the stressed portfolio value.
\end{enumerate}

\subsection{Three Key Historical Scenarios}

\begin{table}[H]
\centering
\begin{tabular}{@{}lp{2.5cm}p{2.5cm}p{2.5cm}@{}}
\toprule
\textbf{Factor} & \textbf{GFC 2008} & \textbf{COVID 2020} & \textbf{SVB 2023} \\
\midrule
S\&P~500 & $-57\%$ (peak-to-trough) & $-34\%$ in 23 days & $-5\%$ (broad); bank stocks $-60\%$ \\
VIX & $14 \to 80$ & $14 \to 82$ & $18 \to 31$ \\
10Y Treasury & $-300$bps & $-120$bps & $+400$bps (2022--23 hiking cycle) \\
BBB Spread & $+400$bps & $+350$bps & $+80$bps (banks $+250$bps) \\
House Prices & $-27\%$ (Case-Shiller, over 2 years) & $+5\%$ (no decline) & flat \\
Unemployment & $5\% \to 10\%$ & $3.5\% \to 14.7\%$ (April 2020) & $3.4\%$ (stable) \\
Duration & Sep 2008 -- Mar 2009 (6 months) & Feb 19 -- Mar 23 (23 days) & Mar 8--12 (3 days for SVB) \\
\bottomrule
\end{tabular}
\caption{Key factor moves during three recent crises.
Each crisis has a distinct ``signature'' --- different factors dominate in each.}
\end{table}

\subsection{Worked Example: Applying the GFC Scenario}

\begin{example}[GFC replay on a current portfolio]\label{ex:gfc-replay}
Consider a \$100M portfolio: 60\% equities, 25\% investment-grade bonds (5Y duration),
15\% high-yield bonds (4Y duration).

Apply the GFC factor moves:
\begin{itemize}
    \item \textbf{Equities} ($-57\%$): $\$60\text{M} \times (-0.57) = -\$34.2\text{M}$.
    \item \textbf{IG bonds}: spread widens by 400bps on a 5Y duration portfolio.
          Price change $\approx -5 \times 4.0\% = -20\%$.
          Loss: $\$25\text{M} \times 0.20 = -\$5.0\text{M}$.
          Partially offset by rates falling 300bps: $+5 \times 3.0\% = +15\%$.
          Gain: $\$25\text{M} \times 0.15 = +\$3.75\text{M}$.
          Net IG loss: $-\$1.25\text{M}$.
    \item \textbf{HY bonds}: spread widens by ~800bps (HY spreads moved more than IG).
          Price change $\approx -4 \times 8.0\% = -32\%$.
          Rate offset: $+4 \times 3.0\% = +12\%$.
          Net: $-20\%$. Loss: $\$15\text{M} \times 0.20 = -\$3.0\text{M}$.
\end{itemize}

\textbf{Total GFC scenario loss}: $\$34.2\text{M} + \$1.25\text{M} + \$3.0\text{M} = \$38.45\text{M}$
(38.5\% of portfolio value).

\textbf{Key insight}: the equity allocation dominates the loss.
The bond portfolio partially benefits from the Treasury rally (flight to quality),
which offsets some of the spread widening on IG bonds.
But HY bonds get hit from both sides during a GFC-type event.
\end{example}

\begin{example}[COVID replay]\label{ex:covid-replay}
Same portfolio, COVID March 2020 moves:
\begin{itemize}
    \item \textbf{Equities} ($-34\%$): $-\$20.4\text{M}$.
    \item \textbf{IG bonds}: spread $+350$bps, rates $-120$bps.
          Net spread impact: $-5 \times 3.5\% + 5 \times 1.2\% = -11.5\%$.
          Loss: $\$25\text{M} \times 0.115 = -\$2.875\text{M}$.
    \item \textbf{HY bonds}: spread $+700$bps, rates $-120$bps.
          Net: $-4 \times 7.0\% + 4 \times 1.2\% = -23.2\%$.
          Loss: $\$15\text{M} \times 0.232 = -\$3.48\text{M}$.
\end{itemize}

\textbf{Total COVID scenario loss}: $\$20.4 + \$2.875 + \$3.48 = \$26.76\text{M}$ (26.8\%).
\end{example}

\begin{example}[SVB/rate shock replay]\label{ex:svb-replay}
The SVB 2023 scenario is fundamentally different: it is a \textbf{rate shock}, not a credit shock.
Equities fall modestly, but the damage comes from rising rates
destroying the value of long-duration fixed-income holdings.

Same portfolio, SVB-style moves (rates $+400$bps, equities $-5\%$, bank stocks $-60\%$):
\begin{itemize}
    \item \textbf{Equities} ($-5\%$ broad market): $-\$3.0\text{M}$.
    \item \textbf{IG bonds}: rates $+400$bps, spread $+80$bps.
          Price change: $-5 \times (4.0\% + 0.8\%) = -24\%$.
          Loss: $\$25\text{M} \times 0.24 = -\$6.0\text{M}$.
    \item \textbf{HY bonds}: rates $+400$bps, spread $+80$bps.
          Price change: $-4 \times 4.8\% = -19.2\%$.
          Loss: $\$15\text{M} \times 0.192 = -\$2.88\text{M}$.
\end{itemize}

\textbf{Total SVB scenario loss}: $\$3.0 + \$6.0 + \$2.88 = \$11.88\text{M}$ (11.9\%).

The SVB scenario is less severe in total but exposes a different vulnerability:
the bond portfolio suffers \emph{without} the offsetting rate decline that cushions
spread-driven losses in the GFC and COVID scenarios.
In a rising-rate environment, there is no flight-to-quality offset.
\end{example}

\begin{keyresult}
Each historical crisis has a distinct \textbf{risk factor signature}:
\begin{itemize}
    \item \textbf{GFC 2008}: massive equity decline + credit spread blowout, partially offset by rate decline.
    \item \textbf{COVID 2020}: rapid equity decline + credit spread blowout + rate decline, compressed into 23 days.
    \item \textbf{SVB 2023}: rate shock without credit spread relief, concentrated in bank-exposed sectors.
\end{itemize}
Running all three reveals different portfolio vulnerabilities.
A portfolio that survives the GFC replay may fail the SVB replay, and vice versa.
This is why multiple historical scenarios are essential.
\end{keyresult}

\begin{interpbox}
Notice that the GFC scenario produces the largest absolute loss (\$38.5M),
but the COVID scenario produces the largest loss \textbf{per unit of time}
(\$26.8M in 23 days vs.\ \$38.5M in 6 months).
The speed of the COVID drawdown meant that portfolios had virtually no time to rebalance.
The SVB scenario, while producing the smallest loss overall,
is the most dangerous for portfolios with significant duration exposure and no hedges ---
exactly the position that caused Silicon Valley Bank to fail.
\end{interpbox}

%% ============================================================================
\section{References}
\label{sec:references}
%% ============================================================================

The following texts provide the regulatory, theoretical, and practical foundations
for stress testing in financial risk management.

\begin{enumerate}[label={[\arabic*]}, leftmargin=2cm]
    \item \textbf{Basel Committee on Banking Supervision}\ (2018).
          \emph{Stress testing principles}.
          Bank for International Settlements.
          The definitive regulatory reference for stress testing frameworks.
          Establishes principles for governance, methodology, and use of results.

    \item \textbf{Board of Governors of the Federal Reserve System}\ (2024).
          \emph{2024 Stress Test Methodology and Results}.
          Federal Reserve.
          Annual DFAST documentation including scenario specifications,
          modeling approaches, and aggregate results for covered institutions.

    \item \textbf{Breuer, T.\ \& Csisz\'ar, I.}\ (2013).
          ``Systematic stress tests with entropic plausibility constraints.''
          \emph{Journal of Banking \& Finance}, 37(5), 1552--1559.
          Rigorous formulation of reverse stress testing using information-theoretic plausibility measures.

    \item \textbf{Kupiec, P.H.}\ (2000).
          ``Stress tests and risk capital.''
          \emph{Journal of Risk}, 2(4), 27--39.
          Early framework for integrating stress test results with capital adequacy assessment.

    \item \textbf{Hull, J.C.}\ (2018).
          \emph{Risk Management and Financial Institutions}. 5th ed.
          Wiley. Chapter~27.
          Accessible treatment of stress testing and scenario analysis
          in the context of bank risk management.

    \item \textbf{McNeil, A.J., Frey, R.\ \& Embrechts, P.}\ (2015).
          \emph{Quantitative Risk Management: Concepts, Techniques and Tools}. 2nd ed.
          Princeton University Press.
          Chapters 9 and 14 cover credit risk modeling and stress testing integration.

    \item \textbf{Rebonato, R.}\ (2010).
          \emph{Coherent Stress Testing: A Bayesian Approach to the Analysis of Financial Stress}.
          Wiley.
          Bayesian framework for constructing internally consistent stress scenarios.

    \item \textbf{Glasserman, P., Kang, C.\ \& Kang, W.}\ (2015).
          ``Stress scenario selection by empirical likelihood.''
          \emph{Quantitative Finance}, 15(1), 25--41.
          Systematic approach to selecting worst-case scenarios from historical data.
\end{enumerate}

\vfill
\hrule
\smallskip
\noindent\textit{Andr\'es Gonz\'alez Ortega --- Risk Analyst Project} \hfill \textit{March 2026}

\end{document}
